\documentclass[10pt,a4paper]{amsart}
\usepackage[margin=19mm]{geometry}
\usepackage{amsmath,amssymb,mathtools}
\usepackage[scr=rsfs,cal=euler]{mathalfa}
\usepackage{xfrac}
\usepackage[unicode,hidelinks,bookmarks=false]{hyperref}
\newcommand{\ie}{i.e.\ }
\newcommand{\cf}{cf.\ }
\newcommand{\resp}{respectively\ }
\newcommand{\Subsection}{\subsection}
\providecommand{\nobreakdash}{\nobreak\hbox{-}}
\setlength{\emergencystretch}{2em}
\multlinegap0pt
\usepackage{amssymb}
\usepackage{mathtools}
\usepackage{tikz}
\newtheorem{assumption}{Assumption}
\newtheorem{theorem}{Theorem}
\newtheorem{proposition}{Proposition}
\newtheorem{lemma}{Lemma}
\newcommand{\N}{{\mathbb N}}
\newcommand{\R}{{\mathbb R}}
\newcommand{\Z}{{\mathbb Z}}
\newcommand{\lorenzorrection}[1]{\#1}%{\textcolor{red}{#1}}}
\title{Coupling of the continuum and semiclassical limit. Part I: convergence of eigenvalues}
\author{Matthias Keller, Lorenzo Pettinari, Christiaan J. F. van de Ven}
\date{Source: arXiv:2602.23156v1 (CC BY 4.0)}
\begin{document}
\maketitle

\noindent\emph{Source and licence.} Coupling of the continuum and semiclassical limit. Part I: convergence of eigenvalues. Version arXiv:2602.23156v1 (CC BY 4.0), distributed by arXiv under the Creative Commons Attribution 4.0 International licence, http://creativecommons.org/licenses/by/4.0/, as recorded in the arXiv licence metadata for this exact version and shown on the abstract page https://arxiv.org/abs/2602.23156v1. Full licence notice: \texttt{licenses/arxiv-2602.23156-CC-BY-4.0.txt}.\par\smallskip

\noindent\textit{Editorial scope.} Counted unit: the article's proposition prop:lower\_bound\_general\_potentials (source0.tex lines 979--982). The article's complete original proof of it is delivered with it (source0.tex lines 983--1041, one \texttt{proof} environment of 59 lines). No mathematics is written, repaired or re-derived: the statement and the proof are the article's own lines, in the article's own notation, and no retained line is substituted or re-flowed.  Retained uncounted as the source-local context the proof needs: lines 227--238; lines 278--284; lines 418--421; lines 807--814; lines 903--908.  Nothing was omitted from any retained span, so the package carries no omission record.  No retained line contains a citation, so the wrapper carries no bibliography and no bibliography entry.  No retained line contains a figure environment, an image-inclusion command or drawing code, so no image is delivered and the package declares no asset dependency.  Editorial formatting only, in addition: an AMS \texttt{amsart} preamble that restates the article's own declarations and macros used by the retained lines, namely the environments \texttt{assumption}, \texttt{theorem}, \texttt{proposition}, \texttt{lemma} on separate counters, together with the standard packages those lines need.  The retained environments print in this document as Assumption 1, Theorem 1, Proposition 1, Lemma 1 in order of appearance, whereas the article prints the same items with the numbering of its own full document.  The complete native source of the article as distributed in the arXiv e-print for this exact version is bundled unchanged as \texttt{sources/195377/source0.tex}.
\begin{assumption}\label{assumption}
The following conditions on $V:\R^d\to \R$ are required:
    \begin{enumerate}
	\item[(V1)] $V\in C^{\infty}(\mathbb{R}^d)$ and $V\geq 0$.
	\item[(V2)] $V$ has $1\leq m <+\infty$ zeros, $a_1,\dots a_m$, such that the Hessian at these points  given by    
    \begin{equation*}%\label{eq: Hessians}
	   D^2V(a_l) = ((\partial_{i}\partial_{j}V)(a_l))_{i,j=1}^d
	\end{equation*}
    is strictly positive definite for all $l=1,\ldots,m$, i.e., $a_1,\ldots, a_m$ are non-degenerate minima for $V$.
	\item[(V3)] $V$ is strictly positive at infinity, i.e., there exists a $c >0 $ and $R_{0}>0$ such that $|x|> R_{0}$ implies $V(x)\geq c$.
    \end{enumerate}
\end{assumption}

\begin{theorem}\label{thm:eigenvalue_asymptotics_harmonic_oscillator_main}
    For all $n\in \mathbb{N}_0$ and $\gamma\in (-1,\infty)$, it holds
    \begin{align*}
        \lim_{N\to \infty} \frac{E_n(H_N)}{\lambda_N} = e_n(V^{harm}),
    \end{align*}
    where $\lambda_N=N^{1-\gamma}$.
\end{theorem}

\begin{proposition}[Lower bound]\label{thm:lowerbound}
    For all $n\in \mathbb{N}_0$ and $\kappa>0$, the $n$-th eigenvalue $E_n(\kappa)$ of the operator $\Delta + v_\kappa$ satisfies for  $\kappa$ small 
    $$\liminf_{\kappa\to 0} \frac{E_n(\kappa)}{\kappa^2} \geq   2n+1 .$$
\end{proposition}

\begin{lemma}[IMS localization formula]\label{lem:IMS}
Let $\eta_0,\ldots,\eta_m:\Z\to [0,1]$ be such that $\sum_{j=0}^m\eta_j^2=1$. Then, for any bounded self-adjoint Laplace type operator $L$ on $\ell^2(\Z^d)$ and a potential $V$, we have for $H=L+V$
\[H= \sum_{j=1}^m \eta_j H \eta_j +\frac12\sum_{j=1}^m [\eta_j,[\eta_j,L]].\]
Furthermore, for all $j=1,\ldots,m$
\[\| [\eta_j,[  \eta_j,H]]\|\leq 2\|L\| C^2,\]
where
\[C=\sup_{|x-y|=1} |\eta_j(x)-\eta_j(y)|.\]
\end{lemma}

\begin{lemma}[Estimating the potential error]\label{lem:potential_difference}
    Let $V:\R^d\to [0,\infty)$ be a potential satisfying Assumption~\ref{assumption}~(V1) and~(V2) and $N\in \N_0$. Then, for any normalized $f$ supported in $Q_{r}(Na)$ with $0<r<N$, we have
    \begin{align*}
       \langle f, (H_N-H_{N,l})f\rangle = O(\lambda_N^2(r/N)^3).
    \end{align*}
\end{lemma}

\begin{proposition}\label{prop:lower_bound_general_potentials}
    Let $V:\R^d\to [0,\infty)$ be a potential satisfying Assumption~\ref{assumption} and let $H_N=\frac{N^2}2\Delta + \lambda_N^2 V_N$ with $\lambda_N=N^{1-\gamma}$, $\gamma\in(-1,1)$. Then, for all  $n\in \N$,  the eigenvalues $E_{n}(H_N)$ of the operator $H_{N}$ satisfy
    $$\liminf_{N\to \infty}\frac{E_{n}(H_N)}{\lambda_N} \ge   e_n(V).$$ 
\end{proposition}

\begin{proof}
    \lorenzorrection{ As in the proof of  Proposition~\ref{thm:lowerbound}, the proof works again by induction on the energy level index $n\geq 0$. As above the base case $n=0$ can be obtained as a special case of the proof below without the inductive hypothesis. Suppose we have proved for all $m\leq n-1$ that\begin{equation*}
    E_m(H_N) \geq \lambda_Ne_{m}(V) + o(\lambda_N).
\end{equation*} We shall now prove that the statement holds true for $m=n$.} 

We employ the IMS localization formula from Lemma~\ref{lem:IMS}. 
    To this end, we choose 
      $ \eta :\R\to[0,1]$ be the given as $\eta(y) =0\vee(2-|y|_\infty )\wedge 1$ similar to the proof of Proposition~\ref{thm:lowerbound}. We then define for $l=1,\ldots,m$ and $0<\delta<(1-\gamma)/2$
    \begin{align*}
        \eta_l(x) = \eta\left(\frac{2\lambda_N^{1/2}}{N^{1+\delta}} |x - Na_l|_\infty \right)\quad\text{and}\quad
        \eta_0(x) = \sqrt{1-\sum_{l=1}^m \eta_l(x)^2}.
    \end{align*}
     Thus, $\eta_l$ is supported in the cube $Q_{r}(a_l)$ and equal to $1$ on the smaller cube $Q_{r/2}(a_l)$ with radius $r={\lambda_N^{1/2}}/{N^{1+\delta}} $.    
    Hence, the upper bound  on $\delta$ ensures that $\eta_k$ and $\eta_l$ for $k\neq l$ have eventually disjoint support for $N\to\infty$.
    Furthermore, we estimate the Lipschitz constant $C$ from Lemma~\ref{lem:IMS} for this choice of partition of unity. To this end, we observe that for $|x-y|=1$, we have
    \begin{align*}
        |\eta_l(x)-\eta_l(y)| %&\leq \left|\eta\left(\frac{\lambda_N^{1/2}}{N^{1+\delta}} |x - Na_l|_\infty \right)-\eta\left(\frac{\lambda_N^{1/2}}{N^{1+\delta}} |y - Na_l|_\infty \right)\right| \\
        %&\leq \frac{\lambda_N^{1/2}}{N^{1+\delta}} | |x - Na_l|_\infty - |y - Na_l|_\infty | 
        \leq \frac{2\lambda_N^{1/2}}{N^{1+\delta}}.
    \end{align*}

    
    We now apply the IMS localization formula from Lemma~\ref{lem:IMS} with this choice  to obtain
    \begin{align*}
        H_N = \eta_0 H_N \eta_0+ \sum_{l=1}^m \eta_l H_{N,l} \eta_l +  \sum_{l=1}^m \eta_l (H_N- H_{N,l}) \eta_l+ \frac{N^2}2\sum_{l=0}^m [\eta_l,[\eta_l,\Delta]].
    \end{align*}
We start to estimate the last term. By the Lipschitz estimate above, we have
    \begin{align*}
        N^2\| [\eta_l,[  \eta_l,\Delta]]\|\leq 2\cdot 2N^2\cdot \left(\frac{2\lambda_N^{1/2}}{N^{1+\delta}}\right)^2=  16\frac{\lambda_N}{N^{2\delta}}=o(\lambda_N)
    \end{align*}
    since $\delta>0$. Secondly, we estimate the difference term using Lemma~\ref{lem:potential_difference} which gives
    \begin{align*}
       \langle f, \eta_l (H_N- H_{N,l}) \eta_l f\rangle = O(\lambda_N^2(r/N)^3) = O\left(\lambda_N^2\left(\frac{\lambda_N^{1/2}}{N^{1+\delta}N}\right)^3\right) = o(\lambda_N)
    \end{align*}
    since $\delta<(1-\gamma)/2$ and $r={\lambda_N^{1/2}}/{N^{1+\delta}} $.

   To estimate the first term, we observe that on the support of $\eta_0$, we have for $N$ large enough using $\Delta\geq0$ and the assumption (V3) that $\liminf_{x\to\infty}V(x)\ge c$
    \begin{align*}
        \eta_0 H_N \eta_0 \geq \lambda_N^2 \eta_0^2 c \geq \lambda_N e_n(V) \eta_0^2.
    \end{align*}
    for all $n$ whenever $N$ is large enough.

\lorenzorrection{Now, we estimate the second term $ \sum_{l} \eta_l H_{N,l} \eta_l $. We let a number $r$ such that $e_{n-1}(V) < r < e_{n}(V)$ be given and denote by
$P_{N,n-1,l}$  the orthogonal projection onto the span of the eigenfunctions of $H_{N,l}$  having energy less or equal than $\lambda_Nr$ and $P_{N,0,l}=0$. Thus, by induction hypothesis and the eigenvalue asymptotics of the harmonic oscillator, Theorem~\ref{thm:eigenvalue_asymptotics_harmonic_oscillator_main}, we have that the sum  $\sum_l P_{N,n-1,l}$ has rank at most $n-1$ for $N$ large enough and trivially $ P_{N,0,l} $ has rank $0$ which is needed for the base case.}
\lorenzorrection{Then,  we have
    \begin{align*}
        H_{N,l} &\geq \lambda_N r(I-P_{N,n-1,l}) + P_{N,n-1,l}H_{N,l}P_{N,n-1,l} + o(\lambda_N)\\&= \lambda_N r+ P_{N,n-1,l}(H_{N,l}-\lambda_N r)P_{N,n-1,l} + o(\lambda_N).
    \end{align*}
Now, the operator
\begin{align*}
  Q_{N,n-1}=  \sum_{l=1}^m \eta_l P_{N,n,l}(H_{N,l}-\lambda_N r)P_{N,n,l} \eta_l 
\end{align*}
has rank at most $n-1$ as discussed above.
Thus, taking all estimates together, we obtain
    \begin{align*}
        H_N \geq \lambda_N r + Q_{N,n-1} + o(\lambda_N).
    \end{align*}
    By the min-max principle and  since $r$ can be chosen arbitrarily close to $e_n(V)$, this yields the desired lower bound on the eigenvalues.}    
\end{proof}

\end{document}
